%%%%%%%%%%%%%%%%%%%%%%%%%%%%%%%%%%%%%%%%%%%%%%%%%%%%%%%%%%%%%%%%%%%%%%
% Lớp 11 | Hàm số mũ và hàm số lôgarít | Phép tính luỹ thừa
% Dạng 1D6?1-3: Điều kiện cho luỹ thừa, căn thức
% Số câu: 15  (sinh tự động bởi tach_cau.py; từ nay thầy sửa trực tiếp file này)
%%%%%%%%%%%%%%%%%%%%%%%%%%%%%%%%%%%%%%%%%%%%%%%%%%%%%%%%%%%%%%%%%%%%%%

%%%%%%%%%%%%%%%% TRẮC NGHIỆM 4 PHƯƠNG ÁN %%%%%%%%%%%%%%%%

% nguồn: Ok_TracNghiemLop11_Lan1
\begin{ex}%[1D6N1-3]%[ma:1d6d5269]
	Với $\alpha $ là số thực bất  kì, mệnh đề nào sai?
	\choice 
		{$\sqrt {10^{\alpha}}=\left(\sqrt  {10}\right)^{\alpha}$}
		{$\sqrt {10^{\alpha}}=10^{\frac{\alpha}{2}}$} 
		 {$\left(10^{\alpha}\right)^2=\left(100\right)^{\alpha}$}
		{\True $\left(10^{\alpha}\right)^2=\left(10\right)^{{\alpha}^2}$} 
	 \loigiai{
		Theo định nghĩa và các tính chất của lũy thừa, ta thấy A,  B, C là các mệnh đề đúng.\\
		Xét mệnh đề D: với  $\alpha =1$, ta có: $\left(10^1\right)^2=100\ne \left(10\right)^{1^2}=10$ \\ nên mệnh đề D sai.
	} 
\end{ex}

% nguồn: Ok_TracNghiemLop11_Lan1
\begin{ex}%[1D6N1-3]%[ma:64ad8e46]
	Nếu $\left(a-2\right)^{\tfrac{1}{4}}<\left(a-2\right)^{\tfrac{1}{3}}$ thì khẳng định nào sau đây  là đúng?
	\choice 
		{$2  < a < 3$}
		 {$a > 2$}
		{$a  < 3$}
		{\True $a  > 3$}
	\loigiai{ 
	Ta có: $\dfrac{1}{4}<\dfrac{1}{3}$ và $\left(a-2\right)^{\tfrac{1}{4}}<\left(a-2\right)^{\tfrac{1}{3}}$ nên $a-2 >  1\Leftrightarrow a > 3$.}
\end{ex}

% nguồn: Ok_TracNghiemLop11_Lan1
\begin{ex}%[1D6N1-3]%[ma:1d635a2e]
	Cho  $a$, $b>0$; $\alpha$, $\beta \in \mathbb{R}$. Mệnh đề nào sau đây \textbf{sai}?
	 \choice
		{$\left(a^{\alpha}\right)^{\tfrac{1}{\beta}}=a^{\tfrac{\alpha }{\beta }}$,  $\beta \neq 0$}
		 {$(a\cdot b)^{\alpha }=a^{\alpha}\cdot b^{\alpha}$}
		{$\dfrac{a^{\alpha}}{a^{\beta}}=a^{\alpha-\beta}$} 
		 {\True $a^{\alpha}\cdot b^{\beta}=(ab)^{\alpha +\beta }$}
	\loigiai{ 
		Với $a$, $b>0$ và $\alpha$, $\beta \in \mathbb{R}$ ta có  khẳng định $a^{\alpha}\cdot b^{\beta}=(ab)^{\alpha +\beta }$ sai, các khẳng định còn lại đúng.
	 }
\end{ex}

% nguồn: Ok_TracNghiemLop11_Lan1
\begin{ex}%[1D6N1-3]%[ma:bca648b8]
	Với  $ a$ là số thực dương tuỳ ý, $\sqrt {a^3}$ bằng
	 \choice
		 {$a^{\frac{1}{6}}$}
		 {$a^{\frac{2}{3}}$}
		{$a^6$} 
		 {\True $a^{\frac{3}{2}}$}
	\loigiai{ 
		Ta có $\sqrt  {a^3}={a^{\frac{3}{2}}}$.
	 }
\end{ex}

% nguồn: Ok_TracNghiemLop11_Lan1
\begin{ex}%[1D6H1-3]%[ma:1dfc0b5e]
	Rút gọn  biểu thức $P=x^{\frac{1}{3}}\cdot \sqrt [6]{x}$ với $ x>0$.
	 \choice
		{\True $P=\sqrt  {x}$}
		 {$P=x^{\frac{1}{8}}$}
		{$P=x^{\frac{2}{9}}$} 
		 {$P=x^2$}
	 \loigiai{
		Ta có  $P=x^{\frac{1}{3}}\cdot \sqrt [6]{x}=x^{\frac{1}{3}}\cdot x^{\frac{1}{6}}=x^{\frac{1}{3}+\frac{1}{6}}=x^{\frac{1}{2}}=\sqrt {x}$
	 }
\end{ex}

% nguồn: Ok_TracNghiemLop11_Lan1
\begin{ex}%[1D6H1-3]%[ma:9125bec1]
	 Tập xác định $\mathscr{D}$ của hàm số $y=\left(x^2-x-2\right)^{-3}$ là
	 \choice
		{$\mathscr{D}=\mathbb{R}$} 
		 {$\mathscr{D}=(-\infty;-1) \cup(2;+\infty)$}
		{$\mathscr{D}=(0;+\infty)$} 
		{\True $\mathscr{D}=\mathbb{R} \setminus\{-1;  2\}$}
	 \loigiai{
		Vì $-3$ là số nguyên âm nên hàm  số $y=\left(x^2-x-2\right)^{-3}$ xác định khi và chỉ khi
		$$x^2-x-2 \ne 0 \Leftrightarrow \heva{& x \ne -1 \\ &  x \ne 2.}$$
		Vậy tập xác định của hàm số là  $\mathscr{D}=\mathbb{R} \setminus\{-1; 2\}$.
	} 
\end{ex}

% nguồn: Ok_TracNghiemLop11_Lan1
\begin{ex}%[1D6H1-3]%[ma:b8e04243]
	Giá trị $\sqrt [3]{2021}\cdot \sqrt [5]{2021}$ viết dưới dạng lũy thữa với số  mũ hữu tỷ là
	 \choice
		{$2021^{\frac{2}{5}}$} 
		 {$2021^{\frac{1}{15}}$}
		{\True  $2021^{\frac{8}{15}}$}
		{$2021^{\frac{1}{10}}$} 
	 \loigiai{
		 $\sqrt [3]{2021}\cdot \sqrt [5]{2021}=2021^{\frac{1}{3}}\cdot 2021^{\frac{1}{5}}=2021^{\frac{1}{3}+\frac{1}{5}}=2021^{\frac{8}{15}}$
	} 
\end{ex}

% nguồn: Ok_TracNghiemLop11_Lan1
\begin{ex}%[1D6H1-3]%[ma:39fb53ad]
	Cho  $ a<0,b>0$. Mệnh đề nào sai?
	\choice 
		{$\sqrt [4]{a^4b^8}=| ab^2 |$} 
		 {$\sqrt [4]{a^4b^8}=| a |b^2$}
		{\True $\sqrt  [4]{a^4b^8}=ab^2$}
		{$\sqrt [4]{a^4b^8}=-ab^2$} 
	 \loigiai{
		 Ta có $\sqrt [4]{a^4b^8}=| ab^2 |=-ab^2$ do $ ab^2<0$ với $ a<0,b>0$
	} 
\end{ex}

% nguồn: Ok_TracNghiemLop11_Lan1
\begin{ex}%[1D6H1-3]%[ma:c63248b7]
	Biết rằng $ a-\dfrac{1}{a}=4,a>0$. Giá trị biểu thức $D=\dfrac{a^{\frac{4}{3}}\cdot \left(a^{\frac{-1}{3}}+a^{\frac{2}{3}}\right)}{a^{\frac{1}{4}}\cdot \left(a^{\frac{3}{4}}+a^{\frac{-1}{4}}\right)}$ là một số có dạng  $ m+\sqrt {n}$ với $ m$, $ n$ là hai số tự nhiên. Khẳng định nào đúng?
	 \choice
		{\True  $ m+n=7$}
		{$  m-n=7$}
		{$2m+n=7$} 
		 {$ m+2n=7$}
	\loigiai{ 
		Ta  có:
		\begin{eqnarray*}
			 a-\dfrac{1}{a}&=&4\\
			\Leftrightarrow  a^2-4a-1&=&0\\
			\Leftrightarrow a&=&2+\sqrt 5 \\ 
			\Rightarrow D&=&\dfrac{a^{\frac{4}{3}}\cdot  \left(a^{\frac{-1}{3}}+a^{\frac{2}{3}}\right)}{a^{\frac{1}{4}}\cdot \left(a^{\frac{3}{4}}+a^{\frac{-1}{4}}\right)}\\
			&=&\dfrac{a+a^2}{a+1}\\ 
			 &=&\dfrac{a\left(a+1\right)}{a+1}\\
			&=&a\\ 
			 &=&2+\sqrt 5.
		\end{eqnarray*}
	} 
\end{ex}

% nguồn: Ok_TracNghiemLop11_Lan1
\begin{ex}%[1D6H1-3]%[ma:1d95d754]
	Cho biểu thức $P=\sqrt [3]{x\cdot  \sqrt [4]{x^3\sqrt {x}}}$, với $ x>0 $. Mệnh đề nào đúng?
	\choice 
		{\True  $P={x^{\frac{15}{24}}}$}
		 {$P=x^{\frac{1}{2}}$}
		{$P={x^{\frac{7}{24}}}$} 
		 {$P={x^{\frac{7}{12}}}$}
	 \loigiai{
		Do: $P=\sqrt [3]{x\sqrt [4]{x^3\cdot x^{\frac{1}{2}}}}=\sqrt [3]{x\sqrt [4]{{x^{\frac{7}{2}}}}}=\sqrt  [3]{x\cdot {x^{\frac{7}{8}}}}=\sqrt [3]{{x^{\frac{15}{8}}}}={x^{\frac{15}{24}}}$.
	} 
\end{ex}

% nguồn: Ok_TracNghiemLop11_Lan1
\begin{ex}%[1D6H1-3]%[ma:69cf5520]
	Tập xác định $\mathscr{D}$ của  hàm số $y=(x-3)^{\tfrac{\pi}{2}}$ là
	 \choice
		{$\mathscr{D}=\mathbb{R}  \setminus\{3\}$}
		{$\mathscr{D}=[3;+\infty)$} 
		{$\mathscr{D}=\mathbb{R}$} 
		{\True $\mathscr{D}=(3;+\infty)$} 
	 \loigiai{
		Vì $\dfrac{\pi}{2}$ là số vô tỉ nên hàm số $y=(x-3)^{\tfrac{\pi}{2}}$ xác  định $\Leftrightarrow x - 3 > 0 \Leftrightarrow x>3$.\\
		Vậy tập xác định của hàm số là  $\mathscr{D}=(3;+\infty)$.
	 }
\end{ex}

% nguồn: Ok_TracNghiemLop11_Lan1
\begin{ex}%[1D6H1-3]%[ma:fe94f577]
	Cho  $ a=\dfrac{1}{256}$ và $ b=\dfrac{1}{27}$. Tính $A=a^{\frac{-3}{4}}+b^{\frac{-4}{3}}$.
	 \choice
		 {$23$}
		{$89$} 
		{\True  $145$}
		 {$26$}
	 \loigiai{
		Thay $ a=\dfrac{1}{256}$, $ b=\dfrac{1}{27}$ vào $A=a^{\frac{-3}{4}}+b^{\frac{-4}{3}}$ ta được\\ 
		$A=a^{\frac{-3}{4}}+b^{\frac{-4}{3}}=\left(\dfrac{1}{256}\right)^{\frac{-3}{4}}+\left(\dfrac{1}{27}\right)^{\frac{-4}{3}}=\left(4^{-4}\right)^{\frac{-3}{4}}+\left(3^{-3}\right)^{\frac{-4}{3}}=4^3+3^4=145$. 
	} 
\end{ex}

% nguồn: Ok_TracNghiemLop11_Lan1
\begin{ex}%[1D6V1-3]%[ma:71a628fe]
	Cho $9^x+9^{-x}=23$. Khi đó biểu thức $A=\dfrac{5+3^x+3^{-x}}{1-3^x-3^{-x}}=\dfrac{a}{b}$ với $\dfrac{a}{b}$  là phân số tối giản và $ a,b\in \mathbb{Z}$. Tích $ a\cdot b$ bằng
	 \choice
		 {\True $-10$}
		{$10$} 
		{$-8$} 
		{$8$} 
	 \loigiai{
		 Ta có:
		\begin{eqnarray*}
			9^x+9^{-x}&=&23\\ 
			\Leftrightarrow  \left(3^x+3^{-x}\right)^2&=&25\\
			\Leftrightarrow 3^x+3^{-x}&=&5, \text{\, vì  \,} 3^x+3^{-x}>0,\forall x\in \mathbb{R}\\
			\Rightarrow  A&=&\dfrac{5+3^x+3^{-x}}{1-3^x-3^{-x}}\\
			 &=&\dfrac{5+5}{1-5}=\dfrac{-5}{2}
		\end{eqnarray*}
		Vậy $ a\cdot  b=-10$.
	} 
\end{ex}


%%%%%%%%%%%%%%%% ĐÚNG/SAI %%%%%%%%%%%%%%%%

% nguồn: Ok_dungsai_lop11_Lan1
\begin{ex}%[1D6H1-3]%[ma:c2406d71]
 Với $x$ là  số thực
  \choiceTF
  {Các số $1^{1{,}5}$; $\left(x^2+2\right)^{-1}$;  $\left(\dfrac{1}{x^2+3}\right)^{-2}$ được viết theo thứ tự tăng dần}
  {\True Nếu $\left(x^2+1{,}5\right)^a>\left(x^2+2\right)^a$ thì $a<0$} 
  {\True $(x-1)^{-2 \sqrt{3}}>(x-1)^{-3  \sqrt{2}} \Leftrightarrow x>2$}
  {$\left(x^2+2 x+4\right)^{m^2-3}<\left(x^2+2  x+3\right)^{m^2-3} \Leftrightarrow 0<m<\sqrt{3}$}
  \loigiai{
 \begin{itemchoice}
   \itemch \textbf{Sai}.\\
  Vì $x^2+2  \geq 2>1$ nên $\left(x^2+2\right)^{-1}=\dfrac{1}{x^2+2}<1$;\\
   Vî $x^2+3 \geq 3>1$ nên $\left(\dfrac{1}{x^2+3}\right)^{-2}=\left(x^2+3\right)^2>1$.\\
  Do  đó, $\left(x^2+2\right)^{-1}<1<\left(\dfrac{1}{x^2+3}\right)^{-2}$.
   \itemch \textbf{Đúng}.\\
   Ta có $1<x^2+1{,}5<x^2+2$.\\
  Mà $\left(x^2+1{,}5\right)^a>\left(x^2+2\right)^a  \Rightarrow a<0$.
  \itemch  \textbf{Đúng}.\\
  Ta có $-2 \sqrt{3}>-3 \sqrt{2}$ và $(x-1)^{-2 \sqrt{3}}>(x-1)^{-3 \sqrt{2}} \Leftrightarrow  x-1>1 \Leftrightarrow x>2$
   \itemch \textbf{Sai}.\\
  Ta có $x^2+2 x+4=(x+1)^2+3 \geq 3>1$ và $x^2+2 x+3=(x+1)^2+2 \geq  2>1$.\\
  Mặt khác $x^2+2 x+4>x^2+2  x+3$ và $\left(x^2+2 x+4\right)^{m^2-3}<\left(x^2+2 x+3\right)^{m^2-3}$.\\
  $\Rightarrow m^2-3<0  \Leftrightarrow-\sqrt{3}<m<\sqrt{3}$.
 \end{itemchoice}
 } 
\end{ex}


%%%%%%%%%%%%%%%% TRẢ LỜI NGẮN %%%%%%%%%%%%%%%%

% nguồn: Ok_TraLoiNgan_Lop11_Lan1
\begin{ex}%[1D6V1-3]%[ma:1ea5ea05]
	Có bao nhiêu giá trị nguyên của tham số $m$ để  hàm số $f(x)=\left(2x^2+mx+2\right)^{\tfrac{1}{2}}$ xác định với mọi $x\in \mathbb{R}$?\shortans{7}
	 \loigiai{
		Do $\dfrac{1}{2}\notin \mathbb{Z}$ nên  hàm số đã cho xác định $\Leftrightarrow 2x^2+mx+2>0$.\\
		 Hàm số đã cho xác định với mọi $x\in \mathbb{R}$ khi và chỉ khi
		$$2x^2+mx+2>0, \forall x \in \mathbb{R} \Leftrightarrow \Delta = m^2-16<0 \Leftrightarrow -4<m<4.$$ 
		 Vì $m \in \mathbb{Z}$ nên $m\in \{-3;-2;-1;0;1;2;3\}$. Vậy có $7$ giá trị nguyên của $m$ thỏa yêu cầu bài toán.
	} 
\end{ex}
